\section{Related Work}
\label{sec:related-work}

\subsection{LLM Reasoning}

Numerous research have explored to enhance the logical reasoning capabilities of LLMs. Chain-of-Thought \cite{wei2023chainofthought} is a pioneering work that mirrors human thought processes in a step-by-step way when tackling complex problems. Self-Consistency with CoT \cite{wang2023selfconsistency} samples multiple reasoning path and selects the most consistent answer. Tree-of-Thoughts \cite{yao2024tree} allows LLMs to determine their next course of action by considering various reasoning paths and self-evaluation choices. Graph-of-Thoughts \cite{besta2024graph} further represents the nonlinear task resolution process of LLMs as an arbitrary graph and reasoning on the graph. 
Additionally, \cite{shum2023automatic} involves creating a pool of CoT candidates and selecting the optimal candidate based on certain conditions. However, these methods either use only a single agent or lack communication between agents, which makes them prone to hallucinations or self-perception errors.
Verification \cite{lightman2023lets} and feedback recording are used to enhancement reasoning capabilities. 
STaR \cite{eric2022star} generates multiple chains of thought, from which effective ones are selected.
\cite{zhou2022large} proposes a method for selecting the optimal prompt from the candidate set. 
Skeleton-of-Thought \cite{ning2023skeleton} firstly generates skeleton of answer, followed by the parallel complete of content for each point in the skeleton, thus accelerating answer generation. Table-of-Thoughts \cite{jin2023tab} enhances the accuracy of reasoning through the structured modeling of the reasoning process.

\subsection{Multi-agent Debate}

In multi-agent collaboration, the MAD approach has been demonstrated as an effective orthogonal enhancement in logical reasoning. \cite{liang2023encouraging} proposes a MAD framework that encourages divergent thinking in LLMs, where a judge manages the debate and obtain a final solution. 
\cite{du2023improving} further investigates the impact of the number of agents and rounds of debate on accuracy. \cite{xu2023toward} proposes a multi-agent collaboration strategy that simulates the academic peer review process. \cite{wang2023apollo} integrates a prior knowledge retrieval into the debate process, thereby enhancing reasoning capabilities. \cite{fu2023improving} employs autonomous enhancement of negotiation strategies using a multi-round negotiation game exploration model with two agents. \cite{chan2023chateval} presents various communication strategies and evaluates the effects of these differing approaches. Corex \cite{sun2023corex} employs collaborative methods such as debate, review, and retrieve among multiple agents. \cite{li2024improving} proposes that the utilization of a sparse communication topology in MAD to enhance performance and mitigate computational costs. Recent studies \cite{becker2025stayfocusedproblemdrift} have identified the critical issue of problem drift – a gradual deviation from the original task objective. ConfMAD \cite{lin2025enhancingmultiagentdebateperformance} integrates confidence expression throughout the debate process to improve effectiveness andperformance. 